%% ═══════════════════════════════════════════════════════════════════════════
\section{Lecture 2: Scaling and Approximation}
%% ═══════════════════════════════════════════════════════════════════════════
\date{16/04/2026}
In this lecture, we conclude the examples of order-of-magnitude estimation (specifically, the shooting star problem) and introduce a new, extremely powerful tool: \textbf{Scaling Arguments}. By comparing a physical system either to another known system or to a different state of itself, we can cancel out many unknown factors and gain deep physical insights.

\subsection{Completing the Shooting Star Estimate}

Continuing from the previous lecture, we aim to estimate the mass of a typical ``shooting star'' (meteor).

\ex{Mass of a Shooting Star}{\textbf{The Setup:} When a meteoroid enters Earth's atmosphere, it burns up, converting its kinetic energy into heat and light. We want to estimate its mass $M$ by observing the visible light it emits.

\medskip
\textbf{Observable Properties:}
\begin{enumerate}[label=\textbf{\arabic*.}]
    \item \textbf{Duration:} The event happens very quickly, typically taking $\Delta t \sim 0.5$--$1$ second.
    \item \textbf{Velocity:} As discussed previously, the typical entry velocity is on the scale of Earth's orbital speed around the Sun: $V \sim \SI{30}{km/s} = 3 \times 10^4\,\text{m/s}$.
    \item \textbf{Brightness (Luminosity):} To estimate the luminous power, we can compare it to a familiar object, such as a car headlight seen from a distance. Suppose a shooting star appears as bright as a car's headlight ($P_\text{ref} \sim \SI{10}{W}$ of \emph{pure optical power}) viewed from a distance of $D_\text{ref} \sim \SI{3}{km}$.
    \item \textbf{Distance to Meteor:} The meteor typically burns up in the upper atmosphere, at an altitude $L \sim \SI{50}{km}$ to $\SI{100}{km}$. For our estimate, let us use $L \sim \SI{100}{km} = 10^5\,\text{m}$.
\end{enumerate}

\medskip
\textbf{Energy Conversion Model:}
The initial kinetic energy is $E_k = \frac{1}{2} M V^2$. By the principle of equipartition (assuming roughly equal partition of the dissipated energy into thermal energy of the plasma and radiated light), the energy radiated as light is a fraction $f_\text{light} \sim \frac{1}{2}$ of the total energy.
Thus, the total light energy emitted is:
\[
    E_\text{light} = f_\text{light} \left(\frac{1}{2}M V^2 \right) \approx \frac{1}{4} M V^2
\]
The average optical power emitted during the time $\Delta t$ is:
\[
    P_\text{meteor} = \frac{E_\text{light}}{\Delta t} = \frac{M V^2}{4 \Delta t}
\]

\medskip
\textbf{Flux Comparison:}
We observe the flux (power per unit area). Assuming the light is emitted isotropically, the flux we observe from the meteor at distance $L$ is matched to the flux of the reference headlight at distance $D_\text{ref}$:
\[
    F_\text{obs} = \frac{P_\text{meteor}}{4\pi L^2} \approx \frac{P_\text{ref}}{4\pi D_\text{ref}^2} = F_\text{ref}
\]
Substituting $P_\text{meteor}$:
\[
    \frac{M V^2}{4 \Delta t \cdot 4\pi L^2} \approx \frac{P_\text{ref}}{4\pi D_\text{ref}^2} \implies \frac{M V^2}{4 \Delta t \cdot L^2} \approx \frac{P_\text{ref}}{D_\text{ref}^2}
\]
Solving for $M$:
\[
    M \approx \frac{4 \Delta t \cdot P_\text{ref}}{V^2} \left(\frac{L}{D_\text{ref}}\right)^2
\]

\medskip
\textbf{Calculation:}
Plugging in our rough numbers ($\Delta t = \SI{0.5}{s}$, $P_\text{ref} = \SI{10}{W}$, $V = \SI{3e4}{m/s}$, $L = \SI{100}{km}$, $D_\text{ref} = \SI{3}{km}$):
\[
    \left(\frac{L}{D_\text{ref}}\right)^2 \approx \left(\frac{100}{3}\right)^2 \approx 1000
\]
\[
    M \approx \frac{4 \times 0.5 \times 10}{(3 \times 10^4)^2} \times 1000 = \frac{20}{9 \times 10^8} \times 1000 \approx 2 \times 10^{-5}\,\text{kg} = \SI{20}{mg}.
\]
\begin{remark}
The result is startling: the brilliant streak across the night sky is caused not by a massive rock, but by a tiny grain of dust weighing barely a few milligrams. The enormous optical power derives entirely from the extremely high velocity ($V^2 \propto 10^9$)!
\end{remark}}

\subsection{Principle 4: Scaling Arguments}

When making approximations, a single direct calculation often propagates huge uncertainties (e.g., getting a duration wrong by a factor of 2 leads to a significant error). \textbf{Scaling} allows us to bypass this by computing the \emph{ratio} of related quantities. This relies on referencing an unknown system against a known one, gracefully canceling out parameters we do not know.

\ex{Estimating Gravity on the Moon}{\textbf{Question:} Suppose we want to estimate the acceleration of gravity on the Moon, $g_\text{moon}$, relative to Earth's $g_\oplus \approx \SI{10}{m/s^2}$, without looking up its mass or radius.

From Newton's Law of Universal Gravitation, $g = \frac{GM}{R^2}$. Taking the ratio against Earth:
\[
    \frac{g_\text{moon}}{g_\oplus} = \frac{G M_\text{moon} / R_\text{moon}^2}{G M_\oplus / R_\oplus^2} = \frac{M_\text{moon}}{M_\oplus} \left(\frac{R_\oplus}{R_\text{moon}}\right)^2
\]
Since mass $M \propto \rho R^3$, we can rewrite this ratio in terms of density ($\rho$):
\[
    \frac{g_\text{moon}}{g_\oplus} = \frac{\rho_\text{moon}}{\rho_\oplus} \left(\frac{R_\text{moon}}{R_\oplus}\right)^3 \left(\frac{R_\oplus}{R_\text{moon}}\right)^2 = \frac{\rho_\text{moon}}{\rho_\oplus} \frac{R_\text{moon}}{R_\oplus}
\]
Notice that the gravitational constant $G$ has completely canceled out. We now just need the density ratio and the radius ratio.

\medskip
\textbf{Estimating the Radius Ratio:}
We can measure the angular size of the Moon, $\theta_\text{moon}$. A finger at arm's length subtends roughly $\SI{1.5}{\degree}$, or $\approx 1/40$ radians. The Moon is about one-third of a finger width, yielding:
\[
    \theta_\text{moon} \approx \frac{1}{120}\ \text{rad}.
\]
The radius of the Moon is $R_\text{moon} \approx \frac{1}{2} \theta_\text{moon} D_\text{moon}$, where $D_\text{moon}$ is the distance to the Moon. How do we find $D_\text{moon}$ in terms of Earth's radius $R_\oplus$? We use Kepler's Third Law ($T \propto R^{3/2}$) and compare the Moon to a Low Earth Orbit (LEO) satellite:
\[
    \frac{T_\text{moon}}{T_\text{LEO}} = \left(\frac{D_\text{moon}}{R_\text{LEO}}\right)^{3/2}
\]
An LEO satellite orbits just above Earth's surface ($R_\text{LEO} \approx R_\oplus$) with a period $T_\text{LEO} \approx \SI{90}{min}$. The Moon's orbital period is $\approx 30$ days.
\[
    \left(\frac{D_\text{moon}}{R_\oplus}\right)^{3/2} \approx \frac{30 \times 24 \times 60\,\text{min}}{90\,\text{min}} \approx 480
\]
\[
    \frac{D_\text{moon}}{R_\oplus} \approx 480^{2/3} \approx 60.
\]
Meaning the distance to the Moon is roughly 60 Earth radii.
Now we return to the radius of the Moon:
\[
    R_\text{moon} = \frac{1}{2} \theta_\text{moon} D_\text{moon} = \frac{1}{2} \left(\frac{1}{120}\right) \left(60 R_\oplus\right) = \frac{1}{4} R_\oplus.
\]
So the radius of the Moon is about $1/4$ that of Earth (the true value is $\approx 1/3.7$).

\medskip
\textbf{Determining the Density Ratio:}
If we tentatively assume the Moon has the same density as Earth ($\rho_\text{moon} \approx \rho_\oplus$), our formula gives:
\[
    \frac{g_\text{moon}}{g_\oplus} \approx 1 \times \frac{1}{4} = 0.25
\]
However, we know experimentally that $g_\text{moon} \approx \frac{1}{6} g_\oplus \approx 0.16$.
The discrepancy signals a physical reality: \textbf{The Moon is less dense than Earth}. This happens because Earth has a dense iron core, whereas the Moon is primarily composed of lighter mantle-like material (supporting the giant-impact hypothesis for the Moon's formation). By scaling, we deduced an internal geophysical property purely from simple kinematic observations!}

\ex{Tidal Forces: Sun vs Moon}{\textbf{Question:} Which exerts a stronger tidal force on Earth's oceans: the Sun or the Moon?

The tidal force (or more precisely, the differential gravitational acceleration) on the surface of Earth due to a body of mass $M$ at distance $D$ is given by the difference in gravity across Earth's radius $R_\oplus$:
\[
    a_\text{tidal} \approx \frac{GM}{(D - R_\oplus)^2} - \frac{GM}{D^2} \approx \frac{2GM R_\oplus}{D^3}.
\]
Let us take the ratio of the Sun's tidal force to the Moon's tidal force:
\[
    \frac{a_\text{tidal, Sun}}{a_\text{tidal, Moon}} = \frac{2G M_\text{sun} R_\oplus / D_\text{sun}^3}{2G M_\text{moon} R_\oplus / D_\text{moon}^3} = \frac{M_\text{sun}}{M_\text{moon}} \left(\frac{D_\text{moon}}{D_\text{sun}}\right)^3.
\]
Here we use an elegant cosmic coincidence: \textbf{the angular sizes of the Sun and the Moon in the sky are nearly identical} (which is why total solar eclipses exist).
\[
    \theta_\text{sun} \approx \theta_\text{moon} \implies \frac{R_\text{sun}}{D_\text{sun}} \approx \frac{R_\text{moon}}{D_\text{moon}} \implies \left(\frac{D_\text{moon}}{D_\text{sun}}\right)^3 \approx \left(\frac{R_\text{moon}}{R_\text{sun}}\right)^3
\]
Substituting this back into the tidal ratio equation:
\[
    \frac{a_\text{tidal, Sun}}{a_\text{tidal, Moon}} \approx \frac{M_\text{sun}}{M_\text{moon}} \left(\frac{R_\text{moon}}{R_\text{sun}}\right)^3 = \frac{M_\text{sun} / R_\text{sun}^3}{M_\text{moon} / R_\text{moon}^3} = \frac{\rho_\text{sun}}{\rho_\text{moon}}.
\]
This highlights the immense power of scaling arguments: \textbf{The ratio of their tidal forces depends solely on the ratio of their average densities!}

We know empirically from observing spring and neap tides that the solar tide is roughly half the amplitude of the lunar tide ($a_\text{tidal, Sun} \approx 0.5 \cdot a_\text{tidal, Moon}$).
Consequently, we can conclude that the Sun's average density is about half of the Moon's density:
\[
    \rho_\text{sun} \approx \frac{1}{2}\rho_\text{moon}.
\]
Since the Moon's density is mostly rock ($\rho_\text{moon} \approx \SI{3}{g/cm^3}$), the average density of the Sun must be $\approx \SI{1.5}{g/cm^3}$ (just slightly denser than liquid water).}

\subsection{Preview: Lunar Retreat and Earth's Spin}

Looking ahead, we can analyze the macroscopic dynamics of the Earth-Moon system. Because Earth rotates faster than the Moon orbits (1 day vs 30 days), the tidal bulges raised by the Moon are dragged slightly ahead of the Earth-Moon axis. The gravitational pull of the Moon on this misaligned bulge creates a torque that slows Earth's rotation while simultaneously transferring angular momentum to the Moon, causing its orbit to expand.

Current laser ranging measurements (using retroreflectors left by Apollo missions) indicate the Moon is receding at $v \approx \SI{3.8}{cm/yr}$. If we crudely extrapolate this linearly:
\[
    \text{Time to reach current distance } D_\text{moon} \approx \frac{\SI{380000}{km}}{\SI{3.8}{cm/yr}} = \frac{3.8 \times 10^{10}\,\text{cm}}{\SI{3.8}{cm/yr}} \approx 10^{10}\ \text{years}.
\]
This naive linear estimate yields 10 billion years. However, Earth is only about 4.5 billion years old! The resolution to this paradox is that the tidal interaction is highly non-linear with distance; when the Moon was closer, the tidal forces (which scale as $1/D^3$) were vastly stronger, causing a much faster recession rate. We will formulate the differential equation for this retreat in the next lecture.
